\scanpage{100}
donc $f$ est régulière \unsure{sss|ssi} on a $\alpha\ne-2$. Dans ce cas, la \uline{conique circonscrite} $C(f\deleted{\unsure{\alpha}})$ est non sing.

On notera que c'est le cas en particulier pour $\alpha=+2$, sauf si $k$ est de car. 2. Mais comme $f_0$ est dégénérée, la conique $C(f)$ est une \uline{parabole}.

Dans le cas $\alpha=-2$, on aura
\[
\begin{aligned}
f(x,y)&=(x^2-x)+(y^2-y)+2xy=(x+y)^2-(x+y)\\
&=(x\unsure{+|-}y)(\unsure{\alpha|x}+y-1)=\lambda(x,y)\big[\lambda(x,y)-1\big]
\end{aligned}
\]
i.e. $f$ est le produit de deux fonctions lin affines ayant $\overline{m}$ partie homogène $x+y$. Ainsi $C(f)$ est la réunion des deux droites parallèles\unsure{,} \inserted{(\uline{distinctes})}, d'équations $x+y=0$, $x+y=1$, qui étaient déjà en évidence dans l'étude du cas II (où on avait cependant utilisé les coordonnées
\[
X=x\unsure{-|+}\tfrac14,\quad Y=y-\tfrac14,\quad x+y=X+Y+\tfrac12,
\]
donc les droites précédentes ont comme équations resp.
\[
X+Y=-\tfrac12,\quad X+Y=+\tfrac12\,\unsure{.}\,).
\]
% archive sheet number: 99
